\chapter{Conclusiones y trabajo futuro}\label{chap:conclusiones}

\section{Cumplimiento de los objetivos}
El objetivo general se cumple en su alcance funcional: el sistema evalúa cada evento durante toda
la sesión, mantiene un puntaje de riesgo con continuidad entre sesiones y ejecuta en el equipo una
respuesta graduada, con la decisión centralizada en la API y la aplicación distribuida en los
agentes. La tabla~\ref{tab:cumplimiento-objetivos} detalla cada objetivo específico.

\begin{longtable}{p{5.6cm}p{6.2cm}p{2.2cm}}
\caption{Cumplimiento de los objetivos específicos}\label{tab:cumplimiento-objetivos}\\
\hline
Objetivo & Evidencia & Estado \\
\hline
\endfirsthead
\hline
Objetivo & Evidencia & Estado \\
\hline
\endhead
1. Agente liviano con canal cifrado y autenticado & Recolección de metadatos implementada (RF-01,
RNF-08); canal [[CODIGO: TLS y autenticación mutua]]; consumo [[CODIGO: RNF-02]] & [[CODIGO]] \\
2. Modelo híbrido evaluado contra metas & Modelo entrenado sobre CLUE-LDS; cumple 1 de las 5 metas de
comportamiento y detección, y la de latencia (tabla~\ref{tab:metas-modelo}) [[CODIGO: actualizar si se reentrena]] & Parcial \\
3. Puntaje con continuidad y latencia $\leq$ 500\,ms & Continuidad y decaimiento implementados
(RF-03, RF-04); percentil 95 de 14,9\,ms y percentil 99 de 148,9\,ms en entorno local & Cumplido \\
4. Respuesta graduada y reversible & Catálogo de seis acciones configurables por nivel (RF-07,
RF-08); la ejecución real de las mitigaciones en el equipo se verifica en forma manual
(sección~\ref{sec:pruebas-e2e}); la terminación de procesos no es reversible por limitación del
sistema operativo (RNF-06) & Cumplido con limitación \\
5. Panel con explicación y configuración en tiempo de ejecución & CU-02 implementado (RF-12 a
RF-17, RNF-07); evaluación heurística con cuatro problemas mayores corregidos & Cumplido \\
6. Reentrenamiento aislado y período de aprendizaje & CU-03 y CU-04 implementados; las tres
corridas reales del reentrenamiento (una automática y dos manuales) produjeron modelos candidatos
que no superaron al vigente y no se promovieron (RF-21) & Parcial \\
7. Credenciales y canal protegidos & [[CODIGO: resultado de seguridad]] & [[CODIGO]] \\
8. Análisis de la normativa de datos personales & Marco legal (sección~\ref{sec:marco-legal}) y
análisis del tratamiento (sección~\ref{sec:tratamiento-datos}); quedan pendientes el plazo de
conservación y la supresión completa de los datos de un usuario & Cumplido \\
9. Pruebas con cobertura $\geq$ 85\,\% & [[CODIGO: cantidad de pruebas y cobertura]]; pruebas de extremo
a extremo y evaluación heurística realizadas & [[CODIGO]] \\
\hline
\end{longtable}
{\small Fuente: elaboración propia.\par}

De los 30 requerimientos funcionales, 28 están implementados y 2 en estado parcial, y de los 15 no
funcionales, [[CODIGO: actualizar conteo con RNF-02, RNF-03 y RNF-13]] (sección~\ref{sec:cumplimiento}).

\section{Resultados principales}
El resultado técnico más relevante es la arquitectura: el sistema cierra el ciclo de detección,
decisión y respuesta en el equipo, que es el hueco que el estado del arte identifica en las
soluciones comerciales, y lo hace con una latencia muy por debajo de la que exige una respuesta en
tiempo casi real. La separación entre el punto de decisión y el de aplicación permite cambiar
umbrales y acciones sin modificar los agentes.

El resultado más débil es la capacidad de detección. Sobre escenarios de ataque inyectados en el
historial de 15 usuarios, la detección atribuible nunca supera un tercio de los usuarios en ningún
umbral, y con el umbral vigente del nivel alto llega al 13,3\,\% en el escenario combinado, mientras
que el 16,8\,\% de los días sin ataque termina en nivel alto. Entre las causas, la única verificada
es que casi todos los eventos del conjunto de entrenamiento tienen el país registrado como
desconocido, lo que impide detectar el cambio de país; las demás son inferidas: la granularidad
diaria de las variables de Isolation Forest, que diluye una ráfaga corta de eventos, y la longitud
de las secuencias del LSTM. Estos resultados
se basan en una muestra chica y en escenarios diseñados por el propio equipo, pero alcanzan para
concluir que el modelo, tal como está entrenado, no detecta de forma confiable ataques de corta
duración. [[CODIGO: actualizar si el reentrenamiento cambia los resultados.]]

El análisis económico muestra que el producto puede ser rentable en un escenario de adopción
moderada, con un VAN de USD 35.084 al 25\,\%, pero que su viabilidad depende de la velocidad de
adopción y del costo de adquirir clientes. El marco legal no impide el uso del sistema, pero exige
que la organización que lo adopte informe a sus empleados y formalice la transferencia
internacional de los datos.

\section{Lecciones aprendidas}
\begin{itemize}
  \item Un conjunto de datos público resuelve la falta de datos reales, pero impone su vocabulario
        y sus limitaciones: traducir la telemetría de un equipo Windows al vocabulario de una
        plataforma de archivos compartidos hace que algunos eventos no aporten a las variables del
        modelo, y la falta de etiquetas obliga a construir una evaluación propia.
  \item Sin etiquetas, la evaluación del modelo necesita metas definidas de antemano y escenarios
        controlados; las métricas no supervisadas por sí solas no dicen si el sistema detecta.
  \item Las limitaciones de seguridad del propio sistema (credenciales y canal) son parte del
        producto y no un detalle de despliegue: un sistema de seguridad se evalúa también por
        cómo se protege a sí mismo.
  \item Verificar en el destino real (la base de datos, el correo recibido, el equipo donde se
        ejecuta la acción) evita dar por cumplidas funciones que solo se observan en un registro.
\end{itemize}

\section{Trabajo futuro}
El trabajo futuro se prioriza según el impacto sobre la principal limitación del sistema:
\begin{enumerate}
  \item \textbf{Capacidad de detección}: incorporar variables con granularidad menor al día,
        evaluar secuencias más cortas o solapadas en el LSTM, y calibrar umbrales y pesos con
        escenarios etiquetados más amplios (sección~\ref{sec:interpretacion}), y revisar el componente
        personal en usuarios con historial escaso.
  \item \textbf{Datos propios}: reentrenar el modelo con telemetría de equipos Windows, para
        eliminar la traducción de vocabulario y la variable de país sin información.
  \item \textbf{Explicabilidad del LSTM}: hoy un puntaje impulsado por ese componente no identifica
        qué lo causó.
  \item \textbf{Conservación de datos}: implementar el cierre de la cadena de hashes por período y
        un plazo de conservación configurable (sección~\ref{sec:tratamiento-datos}).
  \item \textbf{Reversión acotada}: registrar las entradas de denegación que crea el agente y
        revertir solo esas.
  \item \textbf{Correcciones del agente}: informar como fallida la suspensión que no se pudo
        registrar y no intentar envíos nuevos con el \emph{circuit breaker} abierto.
  \item \textbf{Pruebas}: agregar pruebas al \emph{frontend} y a los requerimientos sin prueba
        registrada, automatizar la verificación de las mitigaciones y medir la latencia con el
        Redis gestionado y varios agentes concurrentes. [[CODIGO: ajustar según el resultado de la
        cobertura]]
  \item \textbf{Validación comercial}: medir la disposición a pagar y el costo de adquisición a
        través del canal de MSSP, que son las variables que más influyen en la viabilidad.
  \item [[CODIGO: limitaciones de seguridad que queden abiertas después del trabajo del repo de
        código; integración con un proveedor de identidad real.]]
\end{enumerate}
